\chapter[Sermon 39. The Fourth and Fifth Trumpet Is Expounded, of the Opening of the Bottomless Pit, and of Grasshoppers Creeping Out into the Earth.]{The Fourth and Fifth Trumpet Is Expounded, of the Opening of the Bottomless Pit, and of Grasshoppers Creeping Out into the Earth.}
\label{sermon:039}
\markright{Sermon 39. The Fourth and Fifth Trumpet Is Expounded, of the ...}

And the fourth Angel blew, and a third of the sun was struck, and a third of the moon, and a third of the stars, so that a third of them was darkened. And a third of the day was struck, so that it would not shine, and likewise the night. And I beheld, and heard an Angel flying through the midst of Heaven, saying with a loud voice: Woe, woe, woe, to the inhabitants of the Earth, because of the voices yet to come of the trumpet of the three Angels, which were still to blow.

And the fifth angel blew, and I saw a star fall from Heaven to the Earth. And to him was given the key of the bottomless pit. And he opened the bottomless pit, and there arose a smoke from the pit, as it were the smoke of a great furnace. And the sun and the air were darkened by reason of the smoke of the pit. And there came out of the smoke locusts upon the earth, and to them was given power as the scorpions of the Earth have power. And it was said to them that they should not hurt the grass of the Earth, neither any green thing, neither any tree, but only those men which have not the seal in their foreheads, and to them was commanded that they should not kill them, but that they should be vexed five months, and their pain was as the pain that comes of a scorpion when he has stung a man. And in those days shall men seek death, and shall not find it, and shall desire to die, and death shall flee from them.

\index{Augustine, Saint}\index{Repentance}The fourth trumpet declares a harmful and prolonged conflict, which arose in the church concerning the doctrine of Pelagius. This Pelagius taught that the sin of Adam only harmed him, and not mankind, and therefore that in him all men do not die. He asserted that man has free will, so that he may do good. Neither did he believe that he would be free if he needed the help of God; if he had it, he could more easily do good; if he did not have it, he could nevertheless work it by his own virtue and deserve everlasting life. Therefore, he maintained that our victory is not from the help of God, but from free will, and that remission is not given to the penitent after the grace and mercy of God, but after the merit and work of those who, through repentance, are worthy of God’s mercy, as the residue which Saint Augustine recounts in the one hundred and sixteenth Epistle to Boniface, that Pelagius had renounced; which nevertheless in another place he shows that he had taught it and returned to his vomit, as in the register of heresy, heresy number 88, in the second book, chapter 2, to Boniface. The Manicheans, he says, deny that a good man had the beginning of evil from free will. The Pelagians also say that an evil man has free will sufficiently to fulfill a good precept. The Catholic doctrine reproves both of these, and says to them, “God made man upright,” etc. And to them, “If the Son has made you free, you are truly free.” And in the ninth chapter, the same author says that the will of man to evil is free, but that to do good, it must be made free by the grace of God, which contradicts the Pelagians. And where he says that the evil, which was not before, came from him, it is against the Manicheans. Moreover, in the eighth chapter, Pelagius says that the thing which is good may be accomplished sooner if grace helps thereto. By this addition, “more easily,” he truly signifies that he thinks that, although the help of grace is lacking, he can, albeit with more difficulty, perform what is good by free will. Again, the same in the forty-seventh Epistle to Valent. That man, he says, falls into the error of the Pelagians who supposes the grace of God is given for any merit of man, which grace alone makes man free through Jesus Christ our Lord. But again, he who thinks, when the Lord shall come to judgment, that man is not judged after his works which might now, by reason of his age, use the free choice of will, is nevertheless in error. He says essentially the same thing in the second book, chapter 18, concerning the merits and remission of sins.

\index{Rome}With this doctrine of Pelagius—that is to say, darkened (for so John himself explains a little later, saying that a third of them was obscured, and so on)—the third part of the sun, namely Christ, the true sun of righteousness, was struck. For the Pelagians' doctrine denied the grace of Christ, and with human merit trampled underfoot the merit of Christ. By this, also the third part, that is to say a great part, of the Moon, namely the church, is shown to be struck and darkened; moreover, the third part of the stars, I mean preachers and ministers, have not taught with the light that became them. For histories witness that this heresy has severely infected diverse parts of the world, so that even bishops and learned men have followed this noisome error. A synod of bishops was assembled in Palestine in the East, which drove Pelagius to recant. They also disputed sharply against the Pelagian doctrine in Rome, and councils were assembled which condemned it. Synods were assembled in Africa, and after much reasoning, sentence was pronounced against Pelagius. For many were daily taken with this infection. For the doctrine is pleasing, which even today lacks neither maintainers nor defenders. It seems godly, and for the study of virtue needful, to affirm free will and human merit; again, it appears licentious to attribute all things to God’s grace.

He adds that neither the day shone with a third of its light, nor the night with a third of its light. For just as grace could not be fully understood through the doctrine of Pelagius, neither could sin. And Saint Augustine in the second book of original sin, chapters 23 and 24, says that Christian faith properly consists in the account of two men. For by the one we were sold into sin, by the other redeemed from sin; by the one cast down into death, but by the other delivered to life, and so forth. And while all these things are spoken, they are spoken to this end, that we might beware of those heresies.

And thus far we have spoken of the four trumpets and the greatest conflicts in the church. Three trumpets remain, and before them, a brief introduction is set, so that the minds of the listeners might be stirred.

\index{Judgment, Last}And John says how he saw an angel flying through the midst of heaven, and heard him crying: Woe, woe, woe to the inhabitants of the earth, because of the things that would happen to people when the other three trumpets should be blown. Therefore, to each trumpet is joined a woe. We express this very well in Dutch by “woe, woe, woe.” For the Greeks say, and John wrote in Greek, [? \textup{a word lost here}]. And it signifies truly that the times of the former conflicts were sharp, but that those that follow shall be much sharper and crueler. For I told you in another place that this word “woe” encompasses the evils both of this present life and of the life to come, of body as well as of soul. Therefore, the times of Papistry, Mahometry, and of the last judgment shall be most dangerous.

\index{Jerome, Saint}\index{Arethas of Caesarea}\index{Primasius}Note 39.9: The Complutensian copy has an eagle where we read an angel flying through the midst of heaven; perhaps because it found it so in Arethas. Indeed, the common translation, commonly called Saint Jerome’s, has an eagle for an angel. And therefore Primasius reads it so likewise, which seems to have followed the old translation in all things. But the eagle is swift and of most sharp sight, signifying the almighty knowledge of God and expedition unspeakable in doing of things.

\index{Papacy}Note 39.10: The fifth trumpet describes a most terrible battle, which the Pope instigated by introducing errors into the world, indeed by promoting, establishing, and defending them through his ungracious Locusts that consume everything. He will endure until the end of the world. He will discuss him more fully and appropriately in chapters 13 and 14, and following.

\index{Antichrist}The origin of this evil is attributed to the fall of a star. For a star has fallen from heaven to the earth. Stars, as I showed you at the beginning of this book, around the end of the first chapter, represent the state of ministers, or bishops. For just as stars shine in heaven, so bishops, illuminated with heavenly light, ought to shine in the church, both in doctrine and in an honorable life. And they remain in heaven so long as they perform their duty; they fall to the earth when, forgetting heavenly conversation and doctrine, they think upon earthly things, speak of, and pursue honors, pleasures, and such like corruptions. A little later, he will call him an angel, whom he now calls a star. The Church of Rome was notable and pure, commended also by the praise of the Apostle. It had bishops, that is to say, ministers of the church, under Emperor Constantine, around 32, for the most part very well learned, most holy (yet men), and most glorious martyrs of Christ. Again, from Emperor Constantine to Gregory the Great, there are counted bishops or pastors of the church of Rome, around 32, among whom there were not a few diligent, learned, and godly. But there were also among them those who, blinded by the evil of ambition, began to incline more to seek honors and glorious titles than the doctrine of Christ concerning humility and simplicity, and the example of Christ and the apostles. Christ fled when the people would have chosen and made him king. He said that kings should reign, but that apostles and their successors should serve. If kings had offered them realms and riches, they should not have received them. The stories plainly declare what certain bishops of Rome practiced with the churches of Africa, and how they would have ruled over them. Nevertheless, among the later bishops were figures like Pelagius and Gregory, surnamed the Great, who gravely accused the bishops of Constantinople for attempting to establish the church of Constantinople as chief of all others in the world, and the bishop thereof universal. Neither was Gregory ashamed to say expressly that he is the forerunner of Antichrist, whoever would covet the name or title of universal bishop. But Boniface the third of that name, not long after Gregory’s death, moved nothing herewith, required and obtained from Emperor Phocas that the church of Rome might be called and taken as the chief and head of all churches. By this, the bishops of Rome were plucked out of heaven and cast to the earth, utterly beginning to cleave to earthly things, to care for earthly things, even to aspire to the empire and chief rule and government. Here you have what star fell from heaven to the earth.

And to this star (he calls him afterward the Angel of the bottomless pit) or Bishop (I name one, I understand all of that state and succession in that seat) was given the key of the bottomless pit. Christ truly keeps the key of David: as I showed in the second chapter of this book. He gave to the Apostles the keys of the kingdom of heaven, power to open or to shut heaven—that is to say, the ministry of preaching the Gospel, whereby is shown and assuredly promised the forgiveness of sins and eternal life to believers; and the retaining of sins, and certain damnation is threatened to the unbelievers. No godly man doubts but that these keys were given also unto Bishops of Rome; yet every man knows that the later popes would not use them lawfully, but corrupting the Evangelical truth and infecting the lawful ministry, have obtained counterfeit keys. Therefore, to them of the Prince of darkness is given the key of the bottomless pit, namely corrupt and counterfeit doctrine, and not the apostolic, but apostate ministry, whereby, as it were from hell set open, they have brought forth outrageous errors and superstitions, and ungodliness of all sorts. And I suppose it has happened not without God’s providence that Bishops of Rome are called \textit{Clauigers} or key bearers, and wear keys in their arms. But you shall not understand them to be the keys of the kingdom of heaven, but of the bottomless pit rather: for he is a teacher of errors and of all abomination, author moreover of all wars and dissensions, leading them even unto Hell.

God is indeed the source of unending goodness and all truth. This was revealed in Christ by the Apostles through the preaching of the Gospel, and it refreshes with wholesome water all who thirst for eternal salvation. Isaiah refers to this fountain in chapter 55. And Jeremiah in chapter 2. The Lord also speaks of it in the Gospel according to John, in chapters 4 and 7, and in various other passages.

Against this lively fountain of ever-flowing waters, is set the bottomless pit, unsearchable I say because of the malice of Satan, full of ungodliness, abomination, and a kind of lying. From hence blubbereth up into the world by false teachers and ministers of Antichrist whatever error and abomination there is in the world. For Satan, the father of lies, spreads abroad in the world by his instruments whatever darkness there is.

Therefore the star or angel of the bottomless pit—that is, the Pope or Bishop of Rome—opens the bottomless pit with a key, and shortly thereafter ascends with the smoke of the pit. For I have spoken thus far of the beginning of evil; now I shall speak of its progression and manifestation.

The Pope, through his corrupt ministry, opens the way to Hell, and not to Heaven. From Hell ascends a smoke. Smoke in some places of Scripture is a token of God’s presence, wrath, and vengeance, as when a smoke rose in the Temple of Solomon, 1 Kings 8, Isaiah 6:19. In Exodus 19:18, we read that smoke ascended from the mountain as from a furnace. You read in Psalm 18 that smoke went up in God’s wrath, and fire burned before his face. At present, smoke seems to signify hurtful and devilish opinions. Smoke hurts the eyes and prevents clear sight of the truth. So also does perverse doctrine; it darkens the eyes, takes away judgment, and blinds with error. And rightly do they suffer these things from the smoke of God’s wrath and from the lies of deceivable men, who have forsaken the light of the Gospel and the clarity of God’s truth. Under the name of this infernal smoke are contained the opinions and abominable doctrine that the Bishop of Rome, as he is the prelate of the chief church and Apostolic See, is also to be universal pastor and Apostolic, moreover the head of the militant church, the vicar of Christ on earth, whose voice must be heard as much as Christ’s himself; that he has full power in the church, the keys of the kingdom of heaven, and so forth. That he ordains and gives to all churches bishops or pastors, who should govern all other churches according to the prescription of the church of Rome, and so forth.

But how great this smoke is, and how powerfully it is expressed: it ascends, he says, like the smoke of a great furnace. It signifies that the papist opinions and doctrine are thick, or gross, manifold, and apparent; where in truth they are nothing but smoke and vanity, puffed up and vain. But it is of such a power that it darkens the sun and the air. I have told you often that Christ is the sun of righteousness. And we call the air the wholesome doctrine, by which the souls of the faithful are refreshed. Therefore, by the papist doctrine, the sun and the air—that is, Christ and the Gospel—are obscured. Christ is the universal pastor, the high and only Bishop, the head and health of the faithful, who freely forgives sins; which is preached by the Gospel. This doctrine becomes corrupt when the Pope is admitted as head of the church, with the full power of granting indulgences for all sins. Thus is the sun darkened.

However, the evil proceeds further, and establishes itself in the church with much greater effect. For out of the smoke came forth Locusts upon the Earth. For when, through the false persuasion of corrupt doctrine, the eyes of all men were blinded and did not rightly look upon Christ and his only gospel, and all men revered the Pope as the vicar of Christ, the head of the church, and a man apostolic, as it were the mouth of God, and he now made bishops and priests, and nourished, advanced, and established monks and friars: an infinite multitude of the clergy increased most abundantly, in number that could not be numbered. For he himself immediately in the following words, and with a fuller exposition, declares that he speaks nothing of those little worms, the Locusts. For he says, and it was commanded them, that they should not hurt the grass or hay of the earth (and truly, the clergy does not live on hay), neither any green thing, nor any tree, but men only. As though he should say, I speak nothing of grasshoppers such as in times past destroyed Egypt, but I speak of pestilent men, afflicting men with the poison of doctrine. But a little after, they are so described in every point that no man need to doubt that the false clergy is thereby signified. The which thing Primasius also saw, who in his commentaries upon this book said: he puts the authors of evil doctrine. For just as the Locust hurts with its mouth, so do they tear with their preaching, as we read, greedy wolves not sparing the flock, and so forth. Thus says he. There are also other causes why he likened the false clergy to Locusts. If the Locust is alone, it seems most contemptible; so there is nothing more vile than a solitary monk or friar, priest or sophist. But if they swarm together, they are a terror to men, and cannot be driven away with any force; they eat and destroy all. When the prophet Joel would show a great evil to come, he says that the Locusts will come. Sometimes they sing, leap, live at ease and pleasure, to the loss and hindrance of husbandmen. The same things may you see also in the clergy. I speak nothing here of holy priests, that is, lawful ministers of the church, of good men, honest and learned; I speak nothing of the ancient and holy monks, which were burdensome or grievous to no man, and were no preachers, but very laymen, getting their living with their hands, in the church subject with other faithful to the pastors of the church, and so forth. I speak of the unlawful, sluggers, idle bellies, devourers of victuals, but chiefly of false teachers.

\index{Irenaeus}\index{Turks}\index{Zurich}And doubtless the Popes’ clergy is most rightly compared to grasshoppers or caterpillars. For both are innumerable, and they occupy and consume all things. In times past, the ministers of the churches might be numbered, for the number was but small; neither were unprofitable or unnecessary persons nourished of the church goods. There remains a constitution of the emperor Justinian, where among other things, we ordain that there be not at any time in the sacred great church above sixty priests, men deacons, and one hundred subdeacons, ninety readers, and one hundred and ten, nor above twenty-five singers; that the whole number of the clergy of the greater church may consist in four hundred and twenty-five persons, and besides one hundred doorkeepers, as they term them. Therefore, in the most holy great church of this our noble City of Zurich, and in those three churches to the same united (to wit, in the church of our Lady, Saint Theodore, and Saint Irenaeus, let there be so great a multitude of the clergy. This some of the ministers of this imperial city and most large church established five hundred and twenty-five persons. But how many at this day may you find at Rome, or in another great city, priests, monks, friars, and nuns? They exceed this number four times and more. And to leave out many things that might here be brought in, Pope Pius Sabellicus shows in the ninth book of \textit{Aeneidos}, the seventh chapter, that the sect of gray friars was so greatly multiplied throughout the world that then they held and possessed forty provinces, and under every one diverse cloisters and convents (wardens they call the rulers), and exceeded the number of three score thousand men; in so much that the master of the whole order, whom they call general, has been heard many times to offer the pope, preparing an army against the Turks, thirty thousand fighting men of the order of Saint Francis, which should be well able to serve in the wars, and yet be no hindrance or let to their religion or service. And now who is it that knows not how many orders there be of monks and friars? You may therefore account other orders after the rate of the order of Saint Francis, and though you attribute to every one but the one half of that number, to what a sum will it amount? To these if you add the colleges, more and less, throughout so many dioceses, persons, vicars, chaplains, and parish priests, you will grant that not without cause the popish clergy is compared to locusts.

I need not explain with many words how they seize upon and consume everything. It is commonly said that whenever you see any place, fertile and wholesome, wherever you ride or go, you will find it full of clergy and possessed by religious men.

He also speaks explicitly of the power of these Locusts. He presents them as a parable: and he says that power was given to them as scorpions of the earth have. A scorpion is a flattering and, in appearance, a domestic worm, which suddenly strikes with its tail, or rather with the sting of its tail, and so poisons. Therefore, with flattering words, the clergy of Antichrist deceive and wield power through the poison of venomous doctrine. Thus the Apostle also speaks of false teachers in the sixteenth chapter to the Romans. Through fair speech and flattering, they deceive the hearts of the simple. Their power, therefore, is nothing other than evil doctrine, with which, as with the venom of scorpions, they infect simple Christians, but especially those who despise the doctrine of the Gospel.

For it follows a declaration concerning whom these Locusts may harm. There are two kinds of men. One, willingly and knowingly, will perish, and are the open and professed enemies of the holy Gospel; whom, by God’s just judgment, these Scorpiolocusts destroy with their poison. The other is more simple, erring rather through ignorance than through obstinate malice; these are not stung by the Scorpiolocusts, for they have a seal in their foreheads (as spoken of in chapter 7). For the power of this evil is limited, and not without measure. Therefore, the locusts were given this power, that they should not kill—not those wicked who would rather die than live—but these simple ones. They harm them truly, but not to death. And they torment them for five months. And that torment is the trouble of the conscience, which they torment with threatenings, hypocrisy, and wonderful terrors.

There is added for a comfort, five months. The locusts certainly come out in the month of April, and live until September, and when they have lived completely five months, immediately they die. It signifies therefore that those who are consecrated to godliness shall feel these torments for a little while: neither that the deceivers shall always prevail: but that there shall be spaces to rest and breathe in, wherein the godly through the truth may be recovered. For the locusts destroy not, and are seen all the year long. There seems therefore to be a comparison here in this determinate number, that the sense should be: like as the locusts live not longer than from April to September: so doubtless there is a time prefixed to those seducers, and false Papist clergy. Even thus has also the Apostle Paul himself comforted the church: which after he had prophesied that the church should be wonderfully vexed of hypocrites and false teachers, immediately he adds: and like as Jannes and Jambres resisted Moses, so do these resist the truth, men of a mind corrupt, and lewd as concerning the faith: but they shall prevail no longer. For their madness shall be manifest to all men, like as that was of the other. And Primasius says that they are meant here, which although they were entangled with false doctrines, yet having remorse about the end of their life, they receive God’s verity. Again we see, as I warned you in Chapter 7, that all did not perish who were once entangled with the snares of Antichrist. For at the length through the mercy of God they escaped, and sought the grace of God to be given them through Christ, forsaking all superstitions. We see moreover, by reading of histories, how God has at certain times opened the truth by his faithful ministers, through whose preaching the lewdness of the Locusts is interrupted, that men began to smell them out, and to eschew the same; notwithstanding the regenerated, many times have returned, \&c. And likewise other ministers have returned home, \&c.

And furthermore, he declares how great was or is the force of this evil. Their torment, he says, is like the torment of a scorpion when it has struck a man. At first, there is no great pain felt, but little by little it gathers strength, and at last it aches exceedingly. If remedy is had in time, the poison is not deadly; if it is not taken, he who is stung with it dies. To the declaration of this torment, which men feel in their consciences, pertains this that follows: in those days men shall seek death, and so on. And it is a like phrase of speech in a manner as is that same, “mountains fall upon us and cover us,” and so on, whereof I spoke in chapter 6. And it is the voice of one who is sore afflicted and brought in a manner to despair. Doubtless the popish doctrine of merits, of monastic perfection, and of other such like doctrines, have driven many headlong into despair. Hereunto is added that the times of the locusts were most full of sorrows, whereof all histories complain. The life was not pleasant; the locusts did so set men together by the ears among themselves, and so on. And to be brief, they brought men into such a case that they wished to die. The Lord Jesus deliver us from the poison of these locusts.
